\documentclass[12pt,a4paper]{article}

% --- Преамбула ---
\usepackage[utf8]{inputenc}
\usepackage[T1]{fontenc}
\usepackage[russian]{babel}
\usepackage{amsmath,amssymb}
\usepackage{hyperref}
\usepackage[margin=2.5cm]{geometry}
\usepackage{enumitem}
\usepackage{setspace}
\usepackage{booktabs}
\onehalfspacing

\title{\textbf{Второе дополнение к Субквантовой теории}\\
\large Константа масштаба $C = E + m + O$ для частиц Стандартной Модели}
\author{}
\date{}

\begin{document}
\maketitle

\begin{abstract}
В данном дополнении вводится понятие константы масштаба $C = E + m + O$ --- уникального идентификатора частицы, отражающего её 4D-сложность. На основе трёх моделей ортогональности (линейной, логарифмической и упрощённой) вычислены значения $C$ для всех фундаментальных частиц Стандартной Модели. Показано, что $C$ слабо зависит от выбора эталона ортогональности и может служить устойчивой характеристикой масштаба. Для стабильных частиц $C$ лежит в диапазоне $1$--$2000$, для нестабильных --- достигает десятков и сотен тысяч. Для макрообъектов $C$ пропорциональна их массе. Дополнение содержит полные таблицы значений, готовые для экспериментальной проверки.
\end{abstract}

\section{Определение константы масштаба}

\begin{quote}
\textbf{Определение 1 (Константа масштаба).} Для любой частицы (или макрообъекта) определена величина
\begin{equation}
    C = E + m + O
\end{equation}
где $E$ --- энергия покоя частицы ($E = m$ для покоящейся частицы), $m$ --- её масса (в энергетических единицах), $O$ --- ортогональность. $C$ является уникальным идентификатором масштаба частицы и отражает её 4D-сложность.
\end{quote}

Величина $C$ не зависит от выбора системы отсчёта и является инвариантом для частиц одного типа. Для частиц одного семейства (лептоны, барионы, калибровочные бозоны) $C$ изменяется в ограниченном диапазоне, что делает её удобным классификатором.

\section{Модели ортогональности}

Мы используем три модели для вычисления $O$:

\begin{enumerate}[label=\textbf{Модель \Alph*:}]
    \item \textbf{Линейная модель.} $O = 1 - m/m_p$, где $m_p = 938.272$~МэВ --- масса протона. Протон принят за эталон неортогональности ($O_p = 0$), фотон --- за эталон полной ортогональности ($O_\gamma = 1$).
    \item \textbf{Логарифмическая модель.} $O = 1 - \log_{10}(m/m_\gamma) / 6$, где $m_\gamma$ выбрано так, чтобы $O_e = 0.5$ для электрона. Это даёт $m_\gamma = m_e / 1000 \approx 0.000511$~МэВ.
    \item \textbf{Упрощённая модель.} $C = 2m + O$, где $E$ объединена с $m$ (поскольку для покоящейся частицы $E = m$). Это упрощение удобно для быстрой оценки $C$.
\end{enumerate}

\section{Таблицы значений}

\subsection{Модель A: Линейная}

\begin{table}[h]
    \centering
    \caption{Константа масштаба $C = E + m + O$ (линейная модель)}
    \begin{tabular}{lcccc}
        \toprule
        Частица & Масса, МэВ & $O$ & $E=m$ & $C$ \\
        \midrule
        Фотон $\gamma$ & 0 & 1.000000 & 0 & \textbf{1.00} \\
        Нейтрино $\nu$ & $<10^{-7}$ & $\approx 1.000$ & $\approx 0$ & $\approx 1.00$ \\
        Электрон $e$ & 0.511 & 0.999455 & 0.511 & \textbf{1.51} \\
        Мюон $\mu$ & 105.6 & 0.887 & 105.6 & \textbf{211.3} \\
        Тау $\tau$ & 1777 & --0.894 & 1777 & \textbf{3553} \\
        u-кварк & 2.3 & 0.998 & 2.3 & \textbf{5.6} \\
        d-кварк & 4.8 & 0.995 & 4.8 & \textbf{10.6} \\
        c-кварк & 1275 & --0.359 & 1275 & \textbf{2550} \\
        s-кварк & 95 & 0.899 & 95 & \textbf{191} \\
        t-кварк & 173000 & --183.4 & 173000 & \textbf{345817} \\
        b-кварк & 4180 & --3.456 & 4180 & \textbf{8357} \\
        Протон $p$ & 938.272 & 0 (задано) & 938.272 & \textbf{1876.5} \\
        Нейтрон $n$ & 939.6 & --0.0014 & 939.6 & \textbf{1879.2} \\
        W-бозон & 80400 & --84.7 & 80400 & \textbf{160715} \\
        Z-бозон & 91200 & --96.2 & 91200 & \textbf{182304} \\
        Хиггс $H$ & 125000 & --132.2 & 125000 & \textbf{249868} \\
        \bottomrule
    \end{tabular}
    \label{tab:modelA}
\end{table}

\subsection{Модель B: Логарифмическая}

\begin{table}[h]
    \centering
    \caption{Константа масштаба $C = E + m + O$ (логарифмическая модель)}
    \begin{tabular}{lcccc}
        \toprule
        Частица & Масса, МэВ & $O$ (лог) & $E=m$ & $C$ \\
        \midrule
        Фотон $\gamma$ & 0 & 1.000 & 0 & \textbf{1.00} \\
        Нейтрино $\nu$ & $<10^{-7}$ & $\approx 1.000$ & $\approx 0$ & $\approx 1.00$ \\
        Электрон $e$ & 0.511 & 0.500 (задано) & 0.511 & \textbf{1.52} \\
        Мюон $\mu$ & 105.6 & 0.114 & 105.6 & \textbf{211.4} \\
        Тау $\tau$ & 1777 & --0.090 & 1777 & \textbf{3554} \\
        Протон $p$ & 938.272 & 0 (задано) & 938.272 & \textbf{1876.5} \\
        Нейтрон $n$ & 939.6 & --0.044 & 939.6 & \textbf{1879.3} \\
        W-бозон & 80400 & --0.366 & 80400 & \textbf{160799} \\
        Z-бозон & 91200 & --0.375 & 91200 & \textbf{182400} \\
        Хиггс $H$ & 125000 & --0.401 & 125000 & \textbf{250000} \\
        \bottomrule
    \end{tabular}
    \label{tab:modelB}
\end{table}

\subsection{Модель C: Упрощённая}

\begin{table}[h]
    \centering
    \caption{Константа масштаба $C = 2m + O$ (упрощённая модель)}
    \begin{tabular}{lccc}
        \toprule
        Частица & $2m$, МэВ & $O$ & $C = 2m + O$ \\
        \midrule
        Фотон $\gamma$ & 0 & 1.000 & \textbf{1.00} \\
        Электрон $e$ & 1.022 & 0.999455 & \textbf{2.02} \\
        Мюон $\mu$ & 211.2 & 0.887 & \textbf{212.1} \\
        Тау $\tau$ & 3554 & --0.894 & \textbf{3553} \\
        Протон $p$ & 1876.5 & 0 & \textbf{1876.5} \\
        Нейтрон $n$ & 1879.2 & --0.0014 & \textbf{1879.2} \\
        W-бозон & 160800 & --84.7 & \textbf{160715} \\
        Z-бозон & 182400 & --96.2 & \textbf{182304} \\
        Хиггс $H$ & 250000 & --132.2 & \textbf{249868} \\
        \bottomrule
    \end{tabular}
    \label{tab:modelC}
\end{table}

\section{Анализ}

\subsection{Устойчивость $C$}

Значения $C$ для большинства частиц практически не зависят от выбора модели ортогональности. Для электрона разница между Моделью A (1.51) и Моделью B (1.52) составляет менее 1\%. Для тяжёлых частиц (протон, W, Z, Хиггс) значения совпадают в пределах погрешности округления. Это говорит о том, что $C$ является устойчивой характеристикой, не чувствительной к выбору эталона.

\subsection{Диапазоны $C$}

\begin{itemize}
    \item \textbf{Безмассовые частицы (фотон, нейтрино):} $C \approx 1$. Минимальная 4D-сложность.
    \item \textbf{Лёгкие лептоны (электрон):} $C \approx 1.5$--$2$. Очень малая сложность.
    \item \textbf{Тяжёлые лептоны (мюон, тау):} $C \approx 200$--$3500$. Средняя сложность.
    \item \textbf{Барионы (протон, нейтрон):} $C \approx 1880$. Высокая сложность (5 нулей vs 2 у лептонов).
    \item \textbf{Калибровочные бозоны (W, Z):} $C \approx 160\,000$--$182\,000$. Огромная сложность.
    \item \textbf{Хиггс:} $C \approx 250\,000$. Максимальная сложность среди известных частиц.
\end{itemize}

\subsection{$C$ как идентификатор масштаба}

Для макрообъектов $C$ будет пропорциональна их массе (поскольку $O \approx 0$ для макроскопических тел, а $E = m$). Например, для объекта массой 1~кг ($\approx 5.6 \times 10^{29}$~МэВ) $C \approx 2m \approx 1.1 \times 10^{30}$. Это огромное число, но оно отражает 4D-сложность макрообъекта.

\subsection{Связь с топологией}

Значение $C$ коррелирует с числом нулей в собственной моде частицы:

\begin{itemize}
    \item 0 нулей (фотон): $C \approx 1$.
    \item 2 нуля (лептоны): $C \approx 1.5$--$3500$.
    \item 5 нулей (барионы): $C \approx 1880$.
    \item Много нулей (W, Z, Хиггс): $C \approx 10^5$--$10^6$.
\end{itemize}

Чем больше нулей, тем выше 4D-сложность и тем больше $C$.

\section{Предсказания для эксперимента}

\begin{enumerate}[label=\arabic*.]
    \item \textbf{Измерение $O_e$ через рождение пар.} Экспериментальный протокол, описанный в отдельном документе, позволит измерить $O_e$ и тем самым выбрать между Моделью A ($O_e \approx 0.999455$, $C \approx 1.51$) и Моделью B ($O_e = 0.5$, $C \approx 1.52$).
    \item \textbf{Проверка $C$ для мюона.} Если $C_e \approx 1.5$, то $C_\mu \approx 211$. Измерение $C_\mu$ через анализ шума мюона (Часть~6 Субквантовой теории) должно подтвердить это значение.
    \item \textbf{Проверка $C$ для протона.} $C_p \approx 1876.5$ (Модель A и B) или $C_p \approx 1876.5$ (Модель C). Измерение $C_p$ через топологические антенны на протонной основе должно подтвердить это значение.
    \item \textbf{Макрообъекты.} Для макрообъекта массой $M$ предсказание: $C \approx 2M$ (в энергетических единицах). Это можно проверить, измерив ортогональность макрообъекта (через его реакцию на импульс) и вычислив $C$.
\end{enumerate}

\section{Заключение}

Константа масштаба $C = E + m + O$ является устойчивой, не зависящей от выбора эталона характеристикой частицы, отражающей её 4D-сложность. Значения $C$ для всех фундаментальных частиц Стандартной Модели вычислены и представлены в таблицах. Эти значения могут быть проверены экспериментально через измерение ортогональности и анализ шума частиц. Для макрообъектов $C \approx 2m$ (в энергетических единицах).

Дополнение готово для публикации и экспериментальной проверки.

\end{document}